\documentclass[12pt]{article}
\usepackage[letterpaper, margin=1.5cm]{geometry}
\usepackage{amsmath}
\usepackage{amssymb}
\usepackage{empheq}
\setlength\parindent{0pt}

\begin{document}


CPB06 se define mediante las siguientes ecuaciones:
\[ f(\pmb{\sigma}) = B\left[\left(\left|\Sigma_1\right| - k\Sigma_1\right)^a + \left(\left|\Sigma_2\right| - k\Sigma_2\right)^a + \left(\left|\Sigma_3\right| - k\Sigma_3\right)^a \right]^{1/a} \]
\[ \Sigma_i = \textnormal{eigenvalues}\left(\pmb{\Sigma}\right)\]
\[ \pmb{\Sigma} = \mathbb{C} : \pmb{S} \]
\[ \pmb{S} = \pmb{\sigma} - \frac{1}{3}\text{tr}(\pmb{\sigma})\pmb{I} \]
Para calcular la primera derivada de la función, se utiliza la regla de la cadena:
\[ \frac{\partial f}{\partial \pmb{\sigma}} = \sum_{i=1}^{3}\left(\frac{\partial f}{\partial \Sigma_i} \frac{\partial \Sigma_i}{\partial \pmb{\Sigma}}\right) : \frac{\partial \pmb{\Sigma}}{\partial \pmb{S}} : \frac{\partial \pmb{S}}{\partial \pmb{\sigma}} \]
Donde:
\[ \frac{\partial f}{\partial \Sigma_i} = B\left[\left(\left|\Sigma_1\right| - k\Sigma_1\right)^a + \left(\left|\Sigma_2\right| - k\Sigma_2\right)^a + \left(\left|\Sigma_3\right| - k\Sigma_3\right)^{a} \right]^{\frac{1}{a}-1} \cdot     \left(\left|\Sigma_i\right| - k\Sigma_i\right)^{a-1} \cdot \left(\text{sgn}(\Sigma_i) - k\right) \]
\[ \frac{\partial f}{\partial \Sigma_i} = B^a f^{1-a} \cdot \left(\left|\Sigma_i\right| - k\Sigma_i\right)^{a-1} \cdot \left(\text{sgn}(\Sigma_i) - k\right) =
B^a f^{1-a} \cdot \frac{\left(\left|\Sigma_i\right| - k\Sigma_i\right)^{a}}{\Sigma_i} 
\]

\[ \frac{\partial \Sigma_i}{\partial \pmb{\Sigma}} = \frac{\pmb{\Sigma}^2 - \left( I_1 - \Sigma_i\right)\pmb{\Sigma} + \left(\Sigma_i^2 - I_1\Sigma_i+I_2\right)\pmb{I}}{3\Sigma_i^2 - 2I_1\Sigma_i + I_2} \]

Se multiplica estos dos términos:
\[ \frac{\partial f}{\partial \Sigma_i} : \frac{\partial \Sigma_i}{\partial \pmb{\Sigma}} = \frac{B^a f^{1-a} \cdot \left(\left|\Sigma_i\right| - k\Sigma_i\right)^{a}}{\Sigma_i\left(3\Sigma_i^2 - 2I_1\Sigma_i + I_2\right)} \cdot \left(\pmb{\Sigma}^2 - \left( I_1 - \Sigma_i\right)\pmb{\Sigma} + \left(\Sigma_i^2 - I_1\Sigma_i+I_2\right)\pmb{I}\right) \]



\vspace*{10mm}
\[ \frac{\partial \pmb{\Sigma}}{\partial \pmb{S}} = \mathbb{C} \]
\[ \frac{\partial \pmb{S}}{\partial \pmb{\sigma}} = \mathbb{I} - \frac{1}{3}\pmb{I} \otimes \pmb{I} \]
Para la segunda derivada, se aplica nuevamente la regla de la cadena:
\[ \frac{\partial^2 f}{\partial \pmb{\sigma} \partial \pmb{\sigma}} = \frac{\partial}{\partial \pmb{\sigma}}\left(\frac{\partial f}{\partial \pmb{\sigma}}\right) = \frac{\partial}{\partial \pmb{\sigma}}\left(\frac{\partial f}{\partial \Sigma_i} : \frac{\partial \Sigma_i}{\partial \pmb{\Sigma}} : \frac{\partial \pmb{\Sigma}}{\partial \pmb{S}} : \frac{\partial \pmb{S}}{\partial \pmb{\sigma}}\right) \]
Esto implica calcular las derivadas de cada término involucrado, considerando las dependencias entre ellos. La expresión completa para la segunda derivada es compleja y depende de las propiedades específicas del tensor de rigidez \(\mathbb{C}\) y de los autovectores \(\pmb{P}_i\).
\end{document}
